\chapter{Criterii de convergență (cont.). Serii oarecare}

Pe lîngă criteriile prezentate în seminarul anterior, vor mai fi de folos
și altele, pe care le enumerăm mai jos. În continuare, menționăm că
vom lucra cu o serie de forma $ \sum x_n $ și sîntem în ipoteza
$ x_n \in \RR_+, \forall n \in \NN $.

\index{criteriul!Raabe-Duhamel}
\textbf{Criteriul Raabe-Duhamel:} Fie limita următoare:
\[
    \ell = \lim_{n \to \infty} n \left( \dfrac{x_n}{x_{n+1}} - 1 \right).
\]
Atunci:
\begin{itemize}
    \item Dacă $ \ell > 1 $, seria este convergentă;
    \item Dacă $ \ell < 1 $, seria este divergentă;
    \item Dacă $ \ell = 1 $ criteriul nu decide.
\end{itemize}

În multe situații, criteriul Raabe-Duhamel este folositor cînd criteriul
raportului nu decide. Remarcați că, în acest caz, limita de mai sus
este o nedeterminare de forma $ \infty \cdot 0 $, care de multe ori
se poate calcula și este diferită de 1.

\vspace{1cm}

\index{criteriul!logaritmic}
\textbf{Criteriul logaritmic:} Fie limita:
\[
  \ell = \ds\lim_{n \to \infty} \dfrac{- \ln x_n}{\ln n}.
\]
\begin{itemize}
    \item Dacă $ \ell > 1 $, seria este convergentă;
    \item Dacă $ \ell < 1 $, seria este divergentă;
    \item Dacă $ \ell = 1 $, criteriul nu decide.
\end{itemize}

\vspace{1cm}

\index{criteriul!condensării}
\textbf{Criteriul condensării:} Dacă $ (x_n) $ este un șir
descrescător și cu termeni pozitivi, atunci seriile $ \sum x_n $ și
$ \sum 2^n x_{2^n} $ au aceeași natură.

\vspace{1cm}

\index{criteriul!integral}
\textbf{Criteriul integral:} Fie $ f : (0, \infty) \to [0, \infty) $ o funcție
descrescătoare și definim șirul:
\[
    a_n = \int_1^n f(t) dt.
\]
Atunci seria $ \sum f(n) $ este convergentă dacă și numai dacă șirul $ (a_n) $ este
convergent.

%%%%%%%%%%%%%%%%%%%%%%%%%%%%%%%%%%%%%%%%%%%%%%%%%%%%%%%%%%%%%%%%%%%%%%
\section{Exerciții}

1. Studiați natura următoarelor serii cu termeni pozitivi, cu termenul
general $ x_n $ dat de:
\begin{enumerate}[(a)]
    \item $ x_n = \dfrac{1}{\ln n} $; \hfill(D, integral/comparație)
    \item $ x_n = \dfrac{1}{n \ln n} $; \hfill(D, integral/condensare)
    \item $ x_n = \dfrac{1 \cdot 3 \cdot 5 \cdots (2n - 1)}{2 \cdot 4 \cdot 6 \cdots 2n} $ \hfill(D, Raabe)
    \item $ x_n = \left( 1 - \dfrac{3 \ln n}{2n} \right)^n $ \hfill(C, logaritmic)
    \item $ x_n = \dfrac{1}{n \sqrt[n]{n}} $ \hfill(D, comparație/logaritmic)
    \item $ x_n = \left( \dfrac{1}{\ln n} \right)^{\ln(\ln n)} $; \hfill(D, logaritmic)
    \item $ x_n = n^2 e^{-\sqrt{n}} $ \hfill(C, logaritmic)
    \item $ x_n = \dfrac{a^n (n!)^2}{(2n)!}, a > 0. $ \hfill(raport, discuție $ a ? 4 $)
\end{enumerate}

%%%%%%%%%%%%%%%%%%%%%%%%%%%%%%%%%%%%%%%%%%%%%%%%%%%%%%%%%%%%%%%%%%%%%%
\section{Serii alternante}

\index{serii!alternante}
În continuare, discutăm și cazul cînd termenii seriei pot fi negativi. Dar
vom fi interesați doar de un caz particular, anume acela al seriilor \emph{alternante},
adică acelea în care un termen este negativ, iar celălalt pozitiv. Mai precis,
o serie $ \sum x_n $ se numește \emph{alternantă} dacă $ x_n \cdot x_{n+1} < 0 $, 
pentru orice $ n $.

Singurul criteriu de convergență pe care îl folosim pentru aceste cazuri este:

\index{criteriul!Leibniz}
\textbf{Criteriul lui Leibniz:} Fie $ \sum_n (-1)^n x_n $ o serie
alternantă. Dacă șirul $ (x_n) $ este descrescător și converge către 0,
seria este convergentă.

De asemenea, vom mai fi interesați și de:
\begin{itemize}
    \item \emph{serii absolut convergente}, adică acele serii pentru care și
        seria modulelor, și seria dată sînt convergente;
    \item \emph{serii semiconvergente}, adică acele serii pentru care seria
        inițială este convergentă, dar seria modulelor este divergentă.
\end{itemize}

Evident, cum $ x \leq |x| $, rezultă că orice serie absolut convergentă este
convergentă, dar reciproca nu este adevărată.

Pentru acest caz, avem:

\index{criteriul!Abel-Dirichlet}
\textbf{Criteriul Abel-Dirichlet:} Presupunem că seria $ \sum x_n $ se mai poate
scrie sub forma $ \sum \alpha_n y_n $, unde $ (\alpha_n) $ este un șir monoton
și mărginit (deci convergent). Dacă și seria $ \sum y_n $ este convergentă,
atunci seria inițială $ \sum \alpha_n y_n $ este convergentă.

Formulare alternativă: Dacă $ (\alpha_n) $ este un șir monoton, care tinde
către 0, iar șirul cu termenul general $ Y_n = y_1 + \dots + y_n $ este
mărginit, atunci seria $ \sum \alpha_n y_n $ este convergentă.

De exemplu, studiem seria $ \sum_n \dfrac{(-1)^n}{n} $. Este o serie
alternantă, deci:
\begin{itemize}
    \item seria modulelor este $ \zeta(1) $, care este divergentă;
    \item pentru seria dată, aplicăm criteriul lui Leibniz, cu șirul $ x_n = \dfrac{1}{n} $,
        care este descrescător către 0, deci seria este convergentă.
\end{itemize}

Concluzia este că seria $ \sum \dfrac{(-1)^n}{n} $ este \emph{semiconvergentă}.

\textbf{Exercițiu:} Folosind criteriul Leibniz, studiați natura seriei cu termenul general:
\[
    x_n = (-1)^n \dfrac{\log_a n}{n}, a > 1.
\]

%%%%%%%%%%%%%%%%%%%%%%%%%%%%%%%%%%%%%%%%%%%%%%%%%%%%%%%%%%%%%%%%%%%%%%
\section{Aproximarea sumelor seriilor convergente}

Presupunem că avem o serie convergentă și alternantă. Se poate arăta foarte
simplu că, dacă notăm cu $ S $ suma seriei, iar cu $ s_n $ suma primilor $ n $
termeni, cu $ x_n $ termenul general al seriei, are loc inegalitatea:
\begin{equation}\label{eq:aprox}
    \varepsilon = | S - s_n | \leq x_{n + 1}.
\end{equation}
Cu alte cuvinte, eroarea aproximației are ordinul de mărime al primului
termen neglijat.

Deocamdată, exemplele simple pe care le studiem sînt de forma:

\textbf{Exercițiu:} Să se aproximeze cu o eroare mai mică decît $ \varepsilon $
sumele seriilor definite de termenul general $ x_n $ de mai jos:
\begin{enumerate}[(a)]
    \item $ x_n = \dfrac{(-1)^n}{n!}, \varepsilon = 10^{-3} $;
    \item $ x_n = \dfrac{(-1)^n}{n^3 \sqrt{n}}, \varepsilon = 10^{-2} $.
\end{enumerate}

În ambele cazuri, se folosește inegalitatea din~\eqref{eq:aprox}, de unde se 
scoate $ n $. Se obține $ n = 6 $ pentru primul exercițiu și $ n = 4 $ pentru
al doilea.

Concluzia este că, pentru a obține valoarea sumei seriei cu o precizie de 3, respectiv
2 zecimale, este suficient să considerăm primii 5, respectiv primii 3 termeni ai seriei. 
Eroarea este comparabilă cu primul termen \emph{neglijat} din serie.
%%% Local Variables:
%%% mode: latex
%%% TeX-master: "../m1-20-21"
%%% End:
